% GENERATED FILE. Edit canonical lessons, then run npm run generate:books.
% canonical-source-hash: fnv1a64:e0d85dca6f88aaec
% canonical-lessons: LA-C65-capillus, LA-C65-venter, LA-C65-tergum, LA-C65-frons, LA-C65-barba

\chapter{Hair, Belly, Back, Forehead, Beard}
\label{ch:la-hair-belly-back-forehead-beard}
\hlchaptermodality{65}

\begin{chapteropening}
\textbf{By the end of this chapter:} \emph{I can name hair, the belly, the back, the forehead and a beard.}

Complete the last lesson of chapter 65: I can name hair, the belly, the back, the forehead and a beard.
\end{chapteropening}

\section[capillus, capillī]{capillus, capillī — a hair, the hair}

\label{lesson:LA-C65-capillus}



\begin{quote}
Before the new one: say the Latin for the heart, then the Latin for blood.
\end{quote}

\subsection*{You'll want to know: capillus, capillī}
\textbf{capillus, capillī} — \textquotedblleft{}a hair, the hair\textquotedblright{}. A \emph{capillary} is hair-thin.

\begin{cousinweb}
\textbf{capillus} is linked to \emph{caput}, \textquotedblleft{}head\textquotedblright{}: it is often read as \emph{caput} plus \emph{pilus}, \textquotedblleft{}a single hair\textquotedblright{}, so head-hair (EtymDB 2.1), while Lewis \& Short see a diminutive beside \emph{caput}.
\end{cousinweb}

\begin{grammarlens}[title={What you've built}]
One more word for this chapter.
\end{grammarlens}

\subsection*{Your turn}
\begin{itemize}
\raggedright
  \item \emph{Say it:} \emph{capillus}
  \item \emph{Say it:} \emph{capillus}, once more
  \item \emph{Say it:} say \emph{capillus}
  \item \emph{Recall:} say the Latin for the heart, then the Latin for blood, then say \emph{capillus} again
\end{itemize}

\begin{tcolorbox}[breakable,colback=teal!4,colframe=teal!35!black,title={Before you move on}]
What does capillus, capillī mean? (\textquotedblleft{}A hair, the hair\textquotedblright{}.) Say it once more.
\end{tcolorbox}
\section[venter, ventris]{venter, ventris — the belly}

\label{lesson:LA-C65-venter}



\begin{quote}
Before the new one: say the Latin for blood, then the Latin for a hair.
\end{quote}

\subsection*{You'll want to know: venter, ventris}
\textbf{venter, ventris} — \textquotedblleft{}the belly\textquotedblright{}. A \emph{ventriloquist} speaks from the belly.

\begin{cousinweb}
\textbf{venter} may be inherited, from a Proto-Indo-European \emph{*wendtri-}, \textquotedblleft{}belly\textquotedblright{} (EtymDB 2.1), though its relatives in other languages are few and uncertain.
\end{cousinweb}

\begin{grammarlens}[title={What you've built}]
One more word for this chapter.
\end{grammarlens}

\subsection*{Your turn}
\begin{itemize}
\raggedright
  \item \emph{Say it:} \emph{venter}
  \item \emph{Say it:} \emph{venter}, once more
  \item \emph{Say it:} say \emph{venter}
  \item \emph{Recall:} say the Latin for blood, then the Latin for a hair, then say \emph{venter} again
\end{itemize}

\begin{tcolorbox}[breakable,colback=teal!4,colframe=teal!35!black,title={Before you move on}]
What does venter, ventris mean? (\textquotedblleft{}The belly\textquotedblright{}.) Say it once more.
\end{tcolorbox}
\section[tergum, tergī]{tergum, tergī — the back}

\label{lesson:LA-C65-tergum}



\begin{quote}
Before the new one: say the Latin for a hair, then the Latin for the belly.
\end{quote}

\subsection*{You'll want to know: tergum, tergī}
\textbf{tergum, tergī} — \textquotedblleft{}the back\textquotedblright{}.

\begin{cousinweb}
\textbf{tergum}, \textquotedblleft{}the back\textquotedblright{}, has no sure origin; no proposal is generally accepted. English \emph{tergiversation}, \textquotedblleft{}turning one's back, evasion\textquotedblright{}, is built on it.
\end{cousinweb}

\begin{grammarlens}[title={What you've built}]
One more word for this chapter.
\end{grammarlens}

\subsection*{Your turn}
\begin{itemize}
\raggedright
  \item \emph{Say it:} \emph{tergum}
  \item \emph{Say it:} \emph{tergum}, once more
  \item \emph{Say it:} say \emph{tergum}
  \item \emph{Recall:} say the Latin for a hair, then the Latin for the belly, then say \emph{tergum} again
\end{itemize}

\begin{tcolorbox}[breakable,colback=teal!4,colframe=teal!35!black,title={Before you move on}]
What does tergum, tergī mean? (\textquotedblleft{}The back\textquotedblright{}.) Say it once more.
\end{tcolorbox}
\section[frōns, frontis]{frōns, frontis — the forehead, the front}

\label{lesson:LA-C65-frons}



\begin{quote}
Before the new one: say the Latin for the belly, then the Latin for the back.
\end{quote}

\subsection*{You'll want to know: frōns, frontis}
\textbf{frōns, frontis} — \textquotedblleft{}the forehead, the front\textquotedblright{}. English \emph{front} and \emph{frontier}.

\begin{cousinweb}
\textbf{frōns}, \emph{frontis}, is often traced to Proto-Indo-European \emph{*bhront-}, \textquotedblleft{}forehead\textquotedblright{} (EtymDB 2.1), but the link is not certain. It is a different word from \emph{frōns}, \emph{frondis}, \textquotedblleft{}leafy branches\textquotedblright{}.
\end{cousinweb}

\begin{grammarlens}[title={What you've built}]
One more word for this chapter.
\end{grammarlens}

\subsection*{Your turn}
\begin{itemize}
\raggedright
  \item \emph{Say it:} \emph{frōns}
  \item \emph{Say it:} \emph{frōns}, once more
  \item \emph{Say it:} say \emph{frōns}
  \item \emph{Recall:} say the Latin for the belly, then the Latin for the back, then say \emph{frōns} again
\end{itemize}

\begin{tcolorbox}[breakable,colback=teal!4,colframe=teal!35!black,title={Before you move on}]
What does frōns, frontis mean? (\textquotedblleft{}The forehead, the front\textquotedblright{}.) Say it once more.
\end{tcolorbox}
\section[barba, barbae]{barba, barbae — a beard}

\label{lesson:LA-C65-barba}



\begin{quote}
Before the new one: say the Latin for the back, then the Latin for the forehead.
\end{quote}

\subsection*{You'll want to know: barba, barbae}
\textbf{barba, barbae} — \textquotedblleft{}a beard\textquotedblright{}. A \emph{barber} trims it.

\begin{cousinweb}
\textbf{barba} is inherited, usually traced to Proto-Indo-European \emph{*bhardheh2}, \textquotedblleft{}beard\textquotedblright{} (EtymDB 2.1); English \emph{beard} is a cousin (Lewis \& Short).
\end{cousinweb}

\begin{grammarlens}[title={What you've built}]
One more word for this chapter.
\end{grammarlens}

\subsection*{Your turn}
\begin{itemize}
\raggedright
  \item \emph{Say it:} \emph{barba}
  \item \emph{Say it:} \emph{barba}, once more
  \item \emph{Say it:} say \emph{barba}
  \item \emph{Recall:} say the Latin for the back, then the Latin for the forehead, then say \emph{barba} again
\end{itemize}

\begin{tcolorbox}[breakable,colback=teal!4,colframe=teal!35!black,title={Before you move on}]
What does barba, barbae mean? (\textquotedblleft{}A beard\textquotedblright{}.) Say it once more.
\end{tcolorbox}
